% !PRINTED 177
\weberpage{177}
The solution of the general equations of degrees $2$, $3$ and
$4$ is thus reducible to form problems of only one dimension. We shall see in a later chapter
that the general equation of degree $5$ is reducible to a binary
form problem, and one may now, as an
immediate extension of the problem of reducing the solution of an
equation to pure equations, pose the question:
what is the least number of dimensions of a form problem by which a given equation can be solved?
The problem would then have to be formulated thus:

\gesperrt{From the roots of a given equation, rational functions are to be formed in the smallest possible number, in such a way that they pass into homogeneous linear functions of themselves when the roots are subjected to the permutations of the Galois group of the given equation.}

In this way F. \gesperrt{Klein} has extended the problem of the algebraic
solution of an equation. To open the way to an answer to the question, he has proved for the general equation
of degrees $6$ and $7$ that they can be reduced to quaternary form problems.

The general equation of degree $n$ is, as we have seen,
immediately equivalent to a form problem of $n$ dimensions.
But the question whether the general equation of degree $n$,
when $n$ is greater than 7, corresponds to a form problem of
fewer than $n$ dimensions has not yet been answered; it must probably be answered in the negative $^{1})$.

\begingroup
\renewcommand{\thefootnote}{1)}
\footnotetext[1]{F. \gesperrt{Klein}, Zur Theorie der allgemeinen Gleichungen $6^{\mathrm{ten}}$ und
$7^{\mathrm{ten}}$ Grades; Mathematische Annalen, Bd. XXVIII, S. 18. Cf. also ``The
Evanston Colloquium, Lectures on Mathematics by \gesperrt{Felix Klein}'', Lecture IX, London und New York, Macmillan and Co. (1894).}
\endgroup

\begin{center}
\S.\ 45.

\gesperrt{Influence of relative invariants.}
\end{center}

In the definition of the form problem in \S.\ 43 we have
used only the absolute invariants of the group $S$. Now there
are, as we have seen, also cases in which, besides the ab-
%CONT>>
